\chapter{The Hilbert Transform and the Analytic Signal}

Look at a single note from a guitar string. The waveform wiggles up and down thousands
of times a second, but riding on top of those wiggles is a slower, gentler shape: the
note swells when you pluck it and then fades away. That slow shape is the
\emph{envelope}. Your ear barely hears the individual wiggles at all -- it hears the
\emph{pitch} (how fast the wiggles go) and the \emph{loudness} (how big the envelope is),
and it tracks both as they change in time.

We already have a tool for asking ``what frequencies are present?'' -- the Fourier
transform. But the Fourier transform answers that question for the whole signal at once,
or, with the short-time trick of Chapter~10, for a little window at a time. It does not
hand us, sample by sample, a clean answer to two very physical questions:

\begin{itemize}
  \item \emph{How big is the signal right now?} -- its \textbf{instantaneous amplitude},
        also called the envelope.
  \item \emph{How fast is it oscillating right now?} -- its \textbf{instantaneous
        frequency}.
\end{itemize}

The Hilbert transform is the machine that lets us answer both, point by point in time.
It does this by quietly completing a real signal into a \emph{complex} one -- the
\textbf{analytic signal} -- whose magnitude is the envelope and whose angle, unwound,
gives the phase whose rate of change is the instantaneous frequency. The whole chapter
is really one geometric idea: a real cosine is the shadow of a rotating arrow (a
phasor), and if we can recover the arrow from its shadow, we can read off the arrow's
length (amplitude) and how fast it is turning (frequency) at every instant.

This connects directly to the I/Q signals of Chapter~23 and to the AM/FM demodulation
that underlies radio. It also gives us a robust envelope follower -- exactly the kind of
thing we will want when the wrist tracker has to decide ``is a muscle burst happening
right now, and how strong is it?'' We will also be clear about what envelopes can
and cannot tell us about the hand.

\section{The shadow problem: recovering the arrow from its projection}
\label{ch25:shadow}

Recall the picture that has run through this whole book. A cosine
\begin{equation}
\label{ch25:cos-shadow}
x(t) = A\cos(\omega t + \phi)
\end{equation}
is the real part of a rotating complex phasor:
\begin{equation}
\label{ch25:phasor}
z(t) = A\, e^{\jj(\omega t + \phi)}
     = A\cos(\omega t+\phi) + \jj\,A\sin(\omega t+\phi),
\qquad x(t) = \operatorname{Re}\{z(t)\}.
\end{equation}
The phasor is an arrow of length $A$ spinning at $\omega$ radians per second. The signal
$x(t)$ we actually measure -- a voltage, a microphone sample, an accelerometer reading --
is just the \emph{shadow} of that arrow on the real axis (Figure~\ref{ch25:fig-phasor}).

\begin{figure}[H]
\centering
\begin{tikzpicture}[scale=1.7]
  % circle
  \draw[gray!50] (0,0) circle (1);
  \draw[->,gray!60] (-1.3,0)--(1.4,0) node[right]{Re};
  \draw[->,gray!60] (0,-1.3)--(0,1.4) node[above]{Im};
  % phasor
  \def\ang{52}
  \draw[->,very thick,structurepurple] (0,0)--(\ang:1) node[midway,above left]{$A$};
  \filldraw[structurepurple] (\ang:1) circle (0.6pt);
  % projection onto real axis
  \draw[dashed,specred] (\ang:1)--({cos(\ang)},0);
  \filldraw[specred] ({cos(\ang)},0) circle (0.8pt);
  \node[specred,below] at ({cos(\ang)},-0.05) {$x(t)$};
  % projection onto imag axis
  \draw[dashed,sigblue] (\ang:1)--(0,{sin(\ang)});
  \filldraw[sigblue] (0,{sin(\ang)}) circle (0.8pt);
  \node[sigblue,left] at (-0.04,{sin(\ang)}) {$\hat{x}(t)$};
  % angle
  \draw[->,gray] (0.35,0) arc (0:\ang:0.35);
  \node[gray] at (28:0.52) {$\omega t{+}\phi$};
\end{tikzpicture}
\caption{A real signal $x(t)=A\cos(\omega t+\phi)$ is the projection (shadow) of a
rotating phasor onto the real axis. The \emph{missing} coordinate is the projection onto
the imaginary axis, $A\sin(\omega t+\phi)$. The Hilbert transform manufactures exactly
that missing coordinate, $\hat{x}(t)$, so we can rebuild the whole arrow.}
\label{ch25:fig-phasor}
\end{figure}

Here is the problem. Given only the shadow $x(t)$, can we recover the arrow $z(t)$? A
shadow loses information: many different arrows cast the same shadow at one instant.
What saves us is that the arrow is \emph{turning steadily}. If we know the shadow for
all time, the missing vertical coordinate is determined -- it is the same wave, shifted a
quarter turn. For our single cosine, the vertical coordinate is $A\sin(\omega t+\phi)$,
which is just $A\cos(\omega t + \phi - \pi/2)$: the cosine delayed by $90^\circ$.

\begin{intuitionbox}
A $90^\circ$ phase shift is the whole secret. ``Sine is cosine shifted by a quarter
period.'' If we had a black box that took \emph{any} signal and shifted every frequency
component by exactly $-90^\circ$ -- not by a fixed time, but by a fixed \emph{angle} --
then feeding it a cosine would give us a sine, and we could staple the two together into
the complex arrow. That black box is the Hilbert transform.
\end{intuitionbox}

Why ``a fixed angle, not a fixed time''? Delaying by a fixed \emph{time} $\tau$ shifts a
component at frequency $\omega$ by $\omega\tau$ radians -- more phase shift for higher
frequencies. We want the \emph{same} quarter-turn for every component, regardless of its
frequency. That is a strange, frequency-dependent delay (a quarter of a long period is
many milliseconds; a quarter of a short period is microseconds). No simple time delay can
do it, which is why the Hilbert transform is its own special operator.

\section{Defining the Hilbert transform by what it does to frequencies}
\label{ch25:defn-freq}

The cleanest way to define the Hilbert transform is in the frequency domain, by stating
exactly the phase rotation we just asked for. Let $X(\jj\omega)$ be the Fourier transform
of $x(t)$. The Hilbert transform $\hat{x}(t) = \mathcal{H}\{x\}(t)$ is the signal whose
spectrum is
\begin{equation}
\label{ch25:H-freq}
\widehat{X}(\jj\omega) = -\jj\,\operatorname{sgn}(\omega)\, X(\jj\omega),
\qquad
\operatorname{sgn}(\omega) =
\begin{cases}
+1, & \omega > 0,\\
0, & \omega = 0,\\
-1, & \omega < 0.
\end{cases}
\end{equation}
That is the entire definition. Multiplying by $-\jj$ is a rotation by $-90^\circ$ in the
complex plane; multiplying by $+\jj$ is $+90^\circ$. So
Equation~\eqref{ch25:H-freq} says:

\begin{itemize}
  \item \textbf{positive} frequencies get rotated by $-90^\circ$ (multiplied by $-\jj$),
  \item \textbf{negative} frequencies get rotated by $+90^\circ$ (multiplied by $+\jj$),
  \item the DC term ($\omega=0$) is killed.
\end{itemize}

Figure~\ref{ch25:fig-phase} plots this phase response: flat at $-\pi/2$ for positive
frequencies, flat at $+\pi/2$ for negative ones.

\begin{whybox}
Why opposite signs for positive and negative frequencies? Because a real cosine is built
from \emph{two} counter-rotating phasors,
$\cos\theta = \tfrac12 e^{\jj\theta} + \tfrac12 e^{-\jj\theta}$. The $+\omega$ phasor
spins counterclockwise; the $-\omega$ phasor spins clockwise. To turn $\cos$ into $\sin$
we must push the counterclockwise one back by $90^\circ$ and the clockwise one forward by
$90^\circ$ -- they have to be rotated in \emph{opposite} angular senses so that their real
parts add up to a sine instead of cancelling. The factor $-\jj\operatorname{sgn}(\omega)$
encodes precisely that. Notice its magnitude is $1$ everywhere (except DC): the Hilbert
transform changes only phase, never the amount of any frequency. It is an all-pass
$90^\circ$ phase shifter.
\end{whybox}

\begin{figure}[H]
\centering
\begin{tikzpicture}
\begin{axis}[
  width=12.5cm, height=4.4cm,
  xlabel={$\omega$}, ylabel={phase of $H(\jj\omega)$},
  xmin=-3.4, xmax=3.4, ymin=-2.1, ymax=2.1,
  xtick={0}, xticklabels={$0$},
  ytick={-1.5708,1.5708}, yticklabels={$-\tfrac{\pi}{2}$,$+\tfrac{\pi}{2}$},
  axis lines=middle, clip=false,
]
% phase: -pi/2 for omega>0, +pi/2 for omega<0
\addplot[specred,very thick] coordinates {(-3.3,1.5708) (-0.02,1.5708)};
\addplot[specred,very thick] coordinates {(0.02,-1.5708) (3.3,-1.5708)};
\filldraw[specred] (axis cs:0,0) circle (1.5pt);
\node[specred,anchor=south west] at (axis cs:0.1,-1.5){$-\tfrac{\pi}{2}$ for $\omega>0$};
\node[specred,anchor=north east] at (axis cs:-0.1,1.5){$+\tfrac{\pi}{2}$ for $\omega<0$};
\end{axis}
\end{tikzpicture}
\caption{Phase response of the Hilbert transform, $H(\jj\omega)=-\jj\operatorname{sgn}(\omega)$.
The magnitude is $1$ at every frequency (an all-pass), and the phase is a flat
$\mp90^\circ$: positive frequencies rotate by $-\pi/2$, negative ones by $+\pi/2$, with
DC killed. The same quarter-turn at every frequency is what no fixed time delay can
achieve.}
\label{ch25:fig-phase}
\end{figure}

Let us check it on a cosine. The spectrum of $A\cos(\omega_0 t)$ is two spikes,
$\pi A\,[\delta(\omega-\omega_0) + \delta(\omega+\omega_0)]$. Apply
Equation~\eqref{ch25:H-freq}: the spike at $+\omega_0$ is multiplied by $-\jj$, the spike
at $-\omega_0$ by $+\jj$. Inverse-transforming, that is exactly the spectrum of
$A\sin(\omega_0 t)$. So
\begin{equation}
\label{ch25:cos-to-sin}
\mathcal{H}\{\cos(\omega_0 t)\} = \sin(\omega_0 t),
\qquad
\mathcal{H}\{\sin(\omega_0 t)\} = -\cos(\omega_0 t)
\quad (\omega_0>0).
\end{equation}
Cosine becomes sine; sine becomes minus-cosine. Apply $\mathcal{H}$ twice and you have
rotated everything by $180^\circ$, i.e. negated the signal: $\mathcal{H}^2 = -I$ (on
signals with no DC). This is the operator analogue of $\jj^2=-1$, and it is the first
sign that $\mathcal{H}$ is the right tool for building complex signals out of real ones.

\subsection{The time-domain face: convolution with $1/(\pi t)$}
\label{ch25:time-domain}

Multiplication in frequency is convolution in time
(the convolution theorem, Chapter~11). The inverse Fourier transform of the multiplier
$-\jj\operatorname{sgn}(\omega)$ is the function $1/(\pi t)$. Therefore
\begin{equation}
\label{ch25:H-time}
\hat{x}(t) = (h_{\mathcal H} \conv x)(t)
           = \frac{1}{\pi}\,\text{p.v.}\!\!\int_{-\infty}^{\infty}
             \frac{x(\tau)}{t-\tau}\,\dd\tau,
\qquad h_{\mathcal H}(t) = \frac{1}{\pi t}.
\end{equation}
The ``p.v.'' is the Cauchy \emph{principal value}: the kernel $1/(\pi t)$ blows up at
$t=0$, and we take the symmetric limit around the singularity so the infinity cancels
itself. Geometrically, Equation~\eqref{ch25:H-time} says $\hat x(t)$ is a weighted average
of the whole signal, with weight $1/(t-\tau)$ -- nearby samples pull hardest, with
opposite signs on the two sides of $t$. It is fundamentally \emph{nonlocal}: to know
$\hat x$ at one instant in principle requires the signal at all instants. That is the
price of demanding the same phase shift at every frequency.

\begin{historybox}
David Hilbert (1862--1943) studied this transform around 1905 in the context of
integral equations and what are now called Riemann--Hilbert problems -- pure mathematics,
nothing to do with radios. The engineering payoff came in the 1940s. Dennis Gabor (of
holography fame) introduced the \emph{analytic signal} in his 1946 paper ``Theory of
Communication,'' explicitly to give a precise meaning to the instantaneous amplitude and
frequency of a signal. The mathematician's quarter-turn became the engineer's envelope
detector.
\end{historybox}

\section{The analytic signal: stapling the shadow to its partner}
\label{ch25:analytic}

Now we assemble the complex arrow. Given a real signal $x(t)$, define the
\textbf{analytic signal}
\begin{equation}
\label{ch25:analytic-defn}
x_a(t) = x(t) + \jj\,\hat{x}(t).
\end{equation}
The real part is the original signal (our shadow); the imaginary part is its Hilbert
transform (the manufactured partner). For a single cosine this gives exactly the phasor we
started with:
\begin{equation}
x_a(t) = \cos(\omega_0 t) + \jj\sin(\omega_0 t) = e^{\jj\omega_0 t}.
\end{equation}
We have recovered the rotating arrow from its shadow.

There is a clean frequency-domain way to see what just happened. The spectrum of the
analytic signal is
\begin{equation}
\label{ch25:analytic-spectrum}
X_a(\jj\omega) = X(\jj\omega) + \jj\big(-\jj\operatorname{sgn}(\omega)\big)X(\jj\omega)
              = \big(1 + \operatorname{sgn}(\omega)\big)X(\jj\omega)
              = \begin{cases} 2X(\jj\omega), & \omega>0,\\ X(0), & \omega=0,\\ 0, & \omega<0.\end{cases}
\end{equation}
The analytic signal is the original signal with its \textbf{negative-frequency half deleted}
and its positive half doubled (Figure~\ref{ch25:fig-onesided}).

\begin{figure}[H]
\centering
\begin{tikzpicture}
\begin{axis}[
  name=ax1,
  width=12.5cm, height=4.6cm,
  xlabel={$\omega$}, ylabel={$|X|$},
  xmin=-3.4, xmax=3.4, ymin=0, ymax=2.4,
  xtick={-2,0,2}, xticklabels={$-\omega_0$,$0$,$\omega_0$},
  ytick=\empty, axis lines=left, clip=false,
  every axis plot/.append style={thick},
]
% real signal: two equal spikes
\addplot[sigblue,thick] coordinates {(-2,0) (-2,1)};
\addplot[sigblue,thick] coordinates {(2,0) (2,1)};
\node[sigblue] at (axis cs:-2,1.18){$1$};
\node[sigblue] at (axis cs:2,1.18){$1$};
\node[sigblue,anchor=west] at (axis cs:-3.1,2.05){real $x$: symmetric spectrum};
\end{axis}
\begin{axis}[
  at={(ax1.south west)}, anchor=north west, yshift=-1.4cm,
  width=12.5cm, height=4.6cm,
  xlabel={$\omega$}, ylabel={$|X_a|$},
  xmin=-3.4, xmax=3.4, ymin=0, ymax=2.9,
  xtick={-2,0,2}, xticklabels={$-\omega_0$,$0$,$\omega_0$},
  ytick=\empty, axis lines=left, clip=false,
  every axis plot/.append style={thick},
]
\addplot[specred,thick] coordinates {(2,0) (2,2)};
\node[specred] at (axis cs:2,2.18){$2$};
\node[specred,anchor=west] at (axis cs:-3.1,2.62){analytic $x_a$: negative half deleted, positive half doubled};
\end{axis}
\end{tikzpicture}
\caption{Top: a real cosine has two equal spectral spikes at $\pm\omega_0$. Bottom: its
analytic signal keeps only the positive-frequency spike, doubled. A one-sided spectrum is
what makes the time signal a clean rotating phasor instead of two counter-rotating ones.}
\label{ch25:fig-onesided}
\end{figure}

\begin{happeningbox}
This is the deep reason the analytic signal is so useful. A real signal is forced to have
a \emph{symmetric} spectrum (Hermitian: $X(-\jj\omega)=\overline{X(\jj\omega)}$), because
it is the sum of counter-rotating phasor pairs. Those negative frequencies carry no new
information -- they are just mirror images -- but they are what make a real cosine
``stand still and breathe'' rather than rotate. By throwing the mirror image away, the
analytic signal becomes a signal that purely \emph{rotates}. And a rotating thing has an
unambiguous length (the envelope) and an unambiguous rate of turn (the instantaneous
frequency). You cannot ask ``how fast is $\cos(\omega t)$ rotating?'' -- it does not
rotate, it oscillates. You \emph{can} ask how fast $e^{\jj\omega t}$ rotates. The Hilbert
transform converts the unanswerable question into the answerable one.
\end{happeningbox}

\begin{connectionbox}
This is the same one-sided-spectrum object you met as the \emph{complex baseband} or I/Q
signal in Chapter~23. There, a receiver mixed a real bandpass signal down with two local
oscillators $90^\circ$ apart to produce in-phase ($I$) and quadrature ($Q$) channels;
the complex signal $I + \jj Q$ is exactly an analytic signal (shifted to baseband). The
Hilbert transform is the way to build $I + \jj Q$ from a single real channel by
computation rather than by a second mixer. Linear-algebra reader: deleting negative
frequencies is an orthogonal projection onto the subspace of positive-frequency
exponentials -- the analytic-signal map is a projection followed by a doubling.
\end{connectionbox}

\section{Instantaneous amplitude, phase, and frequency}
\label{ch25:instantaneous}

Once we have the analytic signal in polar form,
\begin{equation}
\label{ch25:polar}
x_a(t) = a(t)\, e^{\jj\theta(t)},
\end{equation}
the two physical quantities fall right out of the geometry of the arrow:
\begin{align}
\label{ch25:env}
a(t) &= \abs{x_a(t)} = \sqrt{x(t)^2 + \hat{x}(t)^2}
       && \text{\textbf{instantaneous amplitude} (envelope),}\\
\label{ch25:phase}
\theta(t) &= \arg x_a(t) = \operatorname{atan2}\big(\hat{x}(t),\,x(t)\big)
       && \text{\textbf{instantaneous phase},}\\
\label{ch25:ifreq}
\omega_i(t) &= \frac{\dd\theta}{\dd t}
       && \text{\textbf{instantaneous (angular) frequency}.}
\end{align}
The envelope $a(t)$ is the length of the arrow -- how big the oscillation is right now,
independent of where it happens to be in its cycle. The instantaneous frequency
$\omega_i(t)$ is how fast the arrow is turning right now. Dividing by $2\pi$ gives the
instantaneous frequency in hertz, $f_i(t) = \omega_i(t)/(2\pi)$.

\begin{example}[Amplitude-modulated tone]
\label{ch25:ex-am}
Let a slowly varying envelope $a(t)>0$ multiply a fast carrier:
\[
x(t) = a(t)\cos(\omega_c t),\qquad \omega_c \text{ large}.
\]
If $a(t)$ varies slowly compared with the carrier (its spectrum sits near DC, well below
$\omega_c$), then the Hilbert transform passes the slow factor through almost untouched
and only shifts the carrier: $\hat{x}(t)\approx a(t)\sin(\omega_c t)$. Then
\[
x_a(t) \approx a(t)\cos(\omega_c t) + \jj\,a(t)\sin(\omega_c t) = a(t)\,e^{\jj\omega_c t},
\]
so the recovered envelope is $|x_a(t)| = a(t)$ exactly, and the instantaneous frequency is
$\dd\theta/\dd t = \omega_c$. The Hilbert transform has \emph{demodulated} the AM signal:
it pulled the message $a(t)$ off the carrier without any tuning, filtering, or knowledge
of $\omega_c$. This is the workhorse use of the analytic signal.
\end{example}

\begin{trapbox}
The clean separation in Example~\ref{ch25:ex-am} requires the \textbf{Bedrosian condition}:
the envelope's spectrum and the carrier's spectrum must not overlap (the envelope must be
\emph{low}-pass and the carrier \emph{high}-pass, with a gap between them). If the
``envelope'' wiggles as fast as the carrier, $\mathcal{H}\{a(t)\cos\omega_c t\}\neq
a(t)\sin\omega_c t$, and the recovered $a(t)=|x_a(t)|$ is \emph{not} the factor you wrote
down. Worse, instantaneous frequency can then go \emph{negative} or spike wildly -- a sign
that ``one envelope and one frequency'' is no longer a sensible description of the signal.
The analytic signal always gives you \emph{an} answer; it is meaningful only when the
signal really is a single narrow-band component (one arrow, slowly changing). For
multi-component signals, decompose first (band-pass filter, or empirical mode
decomposition) and take one component at a time.
\end{trapbox}

\subsection{A worked numeric example: a chirp}
\label{ch25:ex-chirp}

Take a \emph{linear chirp} -- a tone whose frequency sweeps upward -- with a Gaussian
envelope:
\begin{equation}
\label{ch25:chirp}
x(t) = a(t)\cos\!\big(\phi(t)\big),
\qquad
\phi(t) = \omega_0 t + \tfrac{1}{2}\beta t^2,
\qquad
a(t) = e^{-(t/\sigma)^2}.
\end{equation}
The true instantaneous frequency is, by definition of $\phi$,
\[
\omega_i(t) = \frac{\dd\phi}{\dd t} = \omega_0 + \beta t,
\]
a straight line rising at rate $\beta$. Provided the envelope is slow and the sweep is
gentle (Bedrosian again), the analytic signal recovers
$a(t)=|x_a(t)|=e^{-(t/\sigma)^2}$ and $\omega_i(t)=\omega_0+\beta t$ to good
approximation. Figure~\ref{ch25:fig-chirp} shows the chirp with the Hilbert envelope
hugging its peaks, and the recovered instantaneous frequency tracking the rising line.

\begin{figure}[H]
\centering
\begin{tikzpicture}
\begin{axis}[
  width=12.5cm, height=4.6cm,
  xlabel={time $t$}, ylabel={$x(t)$},
  domain=-3:3, samples=200,
  xmin=-3.2, xmax=3.2, ymin=-1.25, ymax=1.45,
  axis lines=left, ytick={-1,0,1}, xtick={-2,0,2},
  legend style={at={(0.02,0.98)},anchor=north west,draw=none,fill=none,font=\footnotesize},
]
% chirp
\addplot[sigblue,thick] {exp(-(x/2.0)^2)*cos(deg(3*x + 0.9*x*x))};
\addlegendentry{$x(t)$ (chirp)}
% envelope +/-
\addplot[specred,thick,dashed] {exp(-(x/2.0)^2)};
\addlegendentry{envelope $a(t)=|x_a(t)|$}
\addplot[specred,thick,dashed,forget plot] {-exp(-(x/2.0)^2)};
\end{axis}
\end{tikzpicture}

\vspace{0.2em}

\begin{tikzpicture}
\begin{axis}[
  width=12.5cm, height=3.8cm,
  xlabel={time $t$}, ylabel={$f_i(t)$},
  domain=-3:3, samples=60,
  xmin=-3.2, xmax=3.2, ymin=0, ymax=1.4,
  axis lines=left, xtick={-2,0,2}, ytick=\empty,
]
\addplot[structurepurple,very thick] {0.48 + 0.14*x};
\node[structurepurple,anchor=south west] at (axis cs:-3.0,1.0)
  {instantaneous frequency rises linearly: $\omega_i=\omega_0+\beta t$};
\end{axis}
\end{tikzpicture}
\caption{A Gaussian-windowed linear chirp. Top: the real signal (blue) with its Hilbert
envelope (red dashed) tracing the amplitude. Bottom: the recovered instantaneous
frequency rises as a straight line -- the analytic signal reads off both ``how loud'' and
``how high'' at every instant, something a single Fourier transform cannot do.}
\label{ch25:fig-chirp}
\end{figure}

\begin{walkbox}
Why differentiate the \emph{unwrapped} phase, and not just take $\operatorname{atan2}$ and
differentiate? Because $\operatorname{atan2}$ returns an angle in $(-\pi,\pi]$ and jumps by
$2\pi$ every time the arrow crosses the negative real axis. Differentiating across a jump
gives a giant spurious spike. We must \emph{unwrap} -- add multiples of $2\pi$ to keep the
phase continuous -- before differentiating. Equivalently, there is a jump-free formula
that never forms the angle at all (next section). Think of it as an odometer: you do not
read instantaneous speed by watching the trip-meter wrap from $999.9$ to $000.0$; you
track total distance and differentiate that.
\end{walkbox}

\subsection{A jump-free instantaneous-frequency formula}
\label{ch25:ifreq-formula}

Differentiating $\theta=\arg x_a = \operatorname{Im}\{\ln x_a\}$ directly gives a tidy
expression that needs no unwrapping. Using $\dd\theta/\dd t=\operatorname{Im}\{\dot
x_a/x_a\}$ and writing $x_a = x + \jj\hat x$,
\begin{equation}
\label{ch25:ifreq-clean}
\omega_i(t)
 = \operatorname{Im}\!\left\{\frac{\dot x_a(t)}{x_a(t)}\right\}
 = \frac{x(t)\,\dot{\hat{x}}(t) - \hat{x}(t)\,\dot{x}(t)}
        {x(t)^2 + \hat{x}(t)^2}.
\end{equation}
The numerator is the cross product of the arrow with its velocity; the denominator is the
arrow's squared length. This is just the formula for angular velocity of a moving point in
the plane -- $\omega = (\vec r \times \dot{\vec r})/\norm{\vec r}^2$ -- specialized to our
arrow. No $\operatorname{atan2}$, no wrapping. It is the recommended way to compute
instantaneous frequency numerically.

\section{Computing it on a real (sampled) signal}
\label{ch25:compute}

In practice $x(t)$ is a finite list of samples, and we want a fast, finite Hilbert
transform. The frequency-domain definition Equation~\eqref{ch25:H-freq} turns directly
into an FFT recipe, which is exactly what \texttt{scipy.signal.hilbert} does -- note that,
confusingly, that function returns the \emph{analytic signal} $x_a$, not $\hat x$ alone.

The recipe is just Equation~\eqref{ch25:analytic-spectrum} on the DFT: take the FFT,
zero the negative-frequency bins, double the positive ones, leave DC (and Nyquist) alone,
and inverse-FFT.

\begin{lstlisting}[language=Python,caption={Analytic signal, envelope, and instantaneous frequency via the FFT.},label={ch25:lst-py}]
import numpy as np

def analytic_signal(x):
    """Return x_a = x + j*Hilbert(x) for a real 1-D array (length N)."""
    N = len(x)
    X = np.fft.fft(x)
    h = np.zeros(N)
    h[0] = 1                      # DC: keep as is
    if N % 2 == 0:
        h[N//2] = 1               # Nyquist bin (even N): keep
        h[1:N//2] = 2             # positive freqs: double
    else:
        h[1:(N+1)//2] = 2         # positive freqs: double
    return np.fft.ifft(X * h)     # negative freqs already zeroed by h

# --- demo: envelope + instantaneous frequency of a chirp ---
fs   = 2000.0                     # sample rate, Hz
t    = np.arange(0, 1.0, 1/fs)
x    = np.exp(-((t-0.5)/0.2)**2) * np.cos(2*np.pi*(50*t + 60*t**2))

xa   = analytic_signal(x)
env  = np.abs(xa)                                  # instantaneous amplitude
phi  = np.unwrap(np.angle(xa))                     # continuous phase
f_i  = np.diff(phi) / (2*np.pi) * fs               # instantaneous freq, Hz
print("freq sweeps from %.0f to %.0f Hz" % (f_i[0], f_i[-1]))
\end{lstlisting}

\begin{trapbox}
Two practical pitfalls. (1) \textbf{Edge effects}: the FFT treats the signal as periodic,
and the nonlocal Hilbert kernel smears the wrap-around discontinuity inward. The envelope
and instantaneous frequency are unreliable near the two ends -- window the signal, or
discard a transient region at each edge. (2) \textbf{Differentiation amplifies noise}:
$\omega_i$ involves a derivative, so even small wiggles in the phase become large spikes
in frequency. Smooth the phase (or the analytic signal) before differentiating, and
remember that instantaneous frequency is meaningless wherever the envelope is near zero
(the arrow's length is tiny, so its angle is undefined and jumps around).
\end{trapbox}

For \emph{streaming} or embedded use you do not want a whole-signal FFT. Instead you
approximate the $90^\circ$ shifter with an FIR filter. The ideal impulse response,
sampled, is
\begin{equation}
\label{ch25:fir}
h[n] = \begin{cases}
\dfrac{2}{\pi n}, & n \text{ odd},\\[1.2ex]
0, & n \text{ even},
\end{cases}
\end{equation}
windowed to a finite length (Figure~\ref{ch25:fig-fir}). Notice every even tap is zero,
which halves the multiplications. A practical odd-length, linear-phase FIR Hilbert filter
delays the signal by $(\text{taps}-1)/2$ samples, so you must delay the real path by the
same amount before combining the two into $x_a$. We design these exactly as the band-pass
filters of Chapter~16, with the Hilbert (Type III/IV) constraint.

\begin{figure}[H]
\centering
\begin{tikzpicture}
\begin{axis}[
  width=12.5cm, height=4.6cm,
  xlabel={tap index $n$}, ylabel={$h[n]$},
  xmin=-8, xmax=8, ymin=-0.8, ymax=0.8,
  xtick={-7,-5,-3,-1,1,3,5,7}, ytick={-0.6,0,0.6},
  axis lines=middle, clip=false,
  ycomb, mark=*, mark size=1.6pt,
]
% h[n] = 2/(pi n), odd n; antisymmetric. 2/pi approx 0.6366
\addplot[sigblue,thick] coordinates {
  (-7,-0.0909) (-5,-0.1273) (-3,-0.2122) (-1,-0.6366)
  (1,0.6366) (3,0.2122) (5,0.1273) (7,0.0909)};
\end{axis}
\end{tikzpicture}
\caption{The ideal sampled Hilbert impulse response $h[n]=2/(\pi n)$ for odd $n$, zero for
even $n$ (only the odd taps are drawn; the even taps sit at zero). It is antisymmetric
about $n=0$ and decays slowly as $1/n$, which is why truncation needs a window. Half the
taps being zero is a free speedup.}
\label{ch25:fig-fir}
\end{figure}

\begin{buildbox}
On an Arduino-class MCU, a full FFT per sample is overkill for envelope following. Two
cheap options:

\textbf{(a) Short FIR Hilbert pair.} Store an odd-length window of recent samples; convolve
with a precomputed $31$- or $63$-tap Hilbert kernel (Equation~\eqref{ch25:fir}) for $\hat
x$, and a matched delay for $x$; then $\text{env}=\sqrt{x^2+\hat x^2}$. With every other
tap zero, a $63$-tap filter is only about $32$ multiply-adds per sample -- fine at audio
rates on an ARM Cortex-M.

\textbf{(b) Rectify-and-smooth.} For pure \emph{envelope} following (not instantaneous
frequency), you often do not need the Hilbert transform at all: take $|x|$ (or $x^2$) and
low-pass filter it with a one-pole leaky integrator,
\texttt{env += alpha*(fabs(x) - env);}. This is the classic analog-style envelope
detector. It is biased (a full-wave-rectified sine has mean $2A/\pi$, not $A$) and ripples
at twice the carrier, but it is one multiply per sample. Use the Hilbert envelope when you
need accuracy or the instantaneous frequency; use rectify-and-smooth when you need
cheapness.
\end{buildbox}

\section{Why this matters: demodulation and envelope following}
\label{ch25:applications}

\begin{applicationbox}
\textbf{AM radio} is literally Example~\ref{ch25:ex-am}: a slow audio message $a(t)$
rides on a fast carrier, and the receiver recovers $a(t)=|x_a(t)|$. Old radios did this
with a diode (rectify-and-smooth); a software-defined radio does it with the Hilbert
transform or with an I/Q mix.

\textbf{FM radio and vibrato} live in the \emph{phase}: the message is encoded as
$\omega_i(t)=\omega_c + k\,m(t)$, so FM demodulation \emph{is} computing instantaneous
frequency, Equation~\eqref{ch25:ifreq-clean}. The same math measures the pitch wobble
(vibrato) of a sung note or the frequency glide of a bird call.

\textbf{Audio and music.} The envelope is the ``A'' in an ADSR amplitude contour; envelope
followers drive compressors, auto-wah, sidechain ducking, and onset detectors. The
instantaneous frequency sharpens pitch estimates between FFT bins (the basis of the
phase-vocoder used for time-stretching and pitch-shifting, Chapter~28).

\textbf{Biomedical.} The Hilbert transform is the standard tool for the \emph{amplitude
envelope} and \emph{phase} of neural oscillations (e.g.\ the power in a band, or
phase--amplitude coupling in EEG/LFP), and for the envelope of the QRS complex in ECG.
\end{applicationbox}

\begin{goalbox}
For the wrist tracker, the analytic signal is our \emph{burst detector}. A surface-EMG or
FMG channel during a muscle contraction looks like a noisy oscillation that swells and
fades -- exactly an amplitude-modulated signal. The Hilbert envelope $|x_a(t)|$ gives a
smooth, sign-independent ``muscle activity level'' without the half-cycle ripple of plain
rectification, which is what we threshold to decide \emph{a gesture is happening now} and
to measure \emph{how hard} (useful for mapping contraction strength to MIDI velocity or to
a continuous controller).

For the IMU, the same trick turns a band-passed accelerometer axis into a tremor or
oscillation envelope -- handy for detecting a periodic gesture (a shake, a tap train) and
reading its strength and rate via instantaneous frequency.

\textbf{Know the limits.} An envelope is a \emph{scalar amplitude}, not a hand
pose. A wrist IMU's Hilbert envelope can tell you ``the forearm is oscillating, fairly
hard, at about 5\,Hz,'' but it cannot recover which finger moved -- a single IMU has no
access to individual finger joints (Chapter~34). Even with EMG, the Hilbert envelope gives
you \emph{when} and \emph{how strong}, not a continuous joint angle; finger-level
discrimination needs multiple EMG/FMG channels and a classifier (Chapters~37--38), and
\emph{continuous} finger pose needs sensors on the fingers themselves. Treat the envelope
as a trigger-and-intensity signal, not as a pose sensor.
\end{goalbox}

\begin{labbox}
Run the code in Listing~\ref{ch25:lst-py} on three signals and watch the envelope and
instantaneous frequency: (1) a pure $440$\,Hz tone -- the envelope should be flat and the
frequency a flat line at $440$; (2) the Gaussian chirp -- envelope a bump, frequency a
ramp; (3) a recording of yourself clenching your fist with a cheap EMG board (or, lacking
one, a band-passed accelerometer axis while you tap a table) -- the envelope should rise
during activity. Then deliberately break the Bedrosian condition: multiply two tones of
similar frequency and watch the ``envelope'' beat and the instantaneous frequency dip
negative. That misbehaviour is the trap of Section~\ref{ch25:instantaneous} made visible.
\end{labbox}

\section{Connection back to I/Q and to projection}
\label{ch25:iq}

We close the loop with Chapter~23. There, a complex baseband signal $I(t)+\jj Q(t)$ came
from \emph{two} physical mixers fed by oscillators $90^\circ$ apart. Here we built the same
kind of object -- a one-sided, purely rotating complex signal -- from a \emph{single} real
channel by computation. The two viewpoints are one thing:

\begin{itemize}
  \item \textbf{Hardware I/Q:} a $90^\circ$ phase split done with analog circuitry, at the
        carrier, before sampling.
  \item \textbf{Hilbert / analytic signal:} the same $90^\circ$ split done with a digital
        filter, after sampling.
\end{itemize}

Both deliver a complex signal whose magnitude is the envelope and whose phase rate is the
instantaneous frequency. Choosing between them is an engineering trade: hardware I/Q needs
two matched analog paths but halves the required sample rate; the Hilbert approach needs
only one ADC channel but must sample fast enough to capture the carrier and then do the
filtering in software.

\begin{formalbox}
\textbf{Summary of the operator.}
$\mathcal{H}$ is the LTI all-pass system with frequency response
$H(\jj\omega)=-\jj\operatorname{sgn}(\omega)$ (unit magnitude, $\mp 90^\circ$ phase) and
impulse response $1/(\pi t)$. It satisfies $\mathcal{H}^2=-I$ (off DC), commutes with time
shifts and derivatives, and is energy-preserving (Parseval): $\norm{\hat x}=\norm{x}$.
The analytic-signal map $x\mapsto x_a=x+\jj\mathcal H x$ has the one-sided spectrum of
Equation~\eqref{ch25:analytic-spectrum}; its modulus is the envelope and its phase
derivative the instantaneous frequency. These are well-defined and physically meaningful
precisely when $x$ is a single narrow-band component (Bedrosian).
\end{formalbox}

\section*{Recap}

\begin{itemize}
  \item A real signal is the shadow of a rotating phasor; the missing coordinate is the
        signal shifted by $90^\circ$ at \emph{every} frequency.
  \item The \textbf{Hilbert transform} supplies that coordinate: in frequency it multiplies
        by $-\jj\operatorname{sgn}(\omega)$ (a $\mp90^\circ$ all-pass rotation); in time it
        convolves with $1/(\pi t)$.
  \item The \textbf{analytic signal} $x_a=x+\jj\hat x$ has a one-sided spectrum (negative
        frequencies deleted, positive doubled) and so purely rotates.
  \item Its modulus is the \textbf{envelope} (instantaneous amplitude); the derivative of
        its unwrapped phase is the \textbf{instantaneous frequency}
        (Equation~\eqref{ch25:ifreq-clean} avoids unwrapping).
  \item These give parameter-free AM demodulation, FM demodulation, envelope following,
        and pitch tracking -- but only for single narrow-band components (Bedrosian); for
        multi-component signals, separate components first.
  \item It is the computational twin of hardware I/Q, and for the wrist tracker it is a
        robust \emph{when/how-strong} burst detector -- never a finger-pose sensor.
\end{itemize}

\section{Exercises}
\label{ch25:exercises}

\begin{exercise}
\label{ch25:ex-cos}
Using only Equation~\eqref{ch25:H-freq}, compute $\mathcal{H}\{\cos(\omega_0 t)\}$ and
$\mathcal{H}\{\sin(\omega_0 t)\}$ for $\omega_0>0$, and confirm
Equation~\eqref{ch25:cos-to-sin}. Then show $\mathcal{H}^2\{\cos(\omega_0 t)\} =
-\cos(\omega_0 t)$, illustrating $\mathcal{H}^2=-I$.
\end{exercise}

\begin{exercise}
\label{ch25:ex-dc}
What is $\mathcal{H}\{c\}$ for a constant signal $x(t)=c$? Use the frequency-domain
definition and explain physically why the answer must be zero. What does this imply about
the envelope of $x(t)=c + \cos(\omega_0 t)$ computed via the analytic signal -- is it
$1$, $c+1$, or something time-varying? (Hint: the DC term is real and unrotated; sketch
the arrow.)
\end{exercise}

\begin{exercise}
\label{ch25:ex-am-env}
For $x(t)=(1+m\cos\omega_m t)\cos\omega_c t$ with $\omega_m\ll\omega_c$ and $0<m<1$ (standard
AM), write down the analytic signal under the Bedrosian approximation and show the
recovered envelope is $1+m\cos\omega_m t$. What goes wrong if $m>1$ (overmodulation)?
\end{exercise}

\begin{exercise}
\label{ch25:ex-ifreq}
Starting from $\theta=\operatorname{Im}\{\ln x_a\}$ and $x_a=x+\jj\hat x$, derive the
jump-free instantaneous-frequency formula Equation~\eqref{ch25:ifreq-clean}. Verify it on
$x(t)=\cos(\omega_0 t)$ (for which $\hat x=\sin\omega_0 t$) and confirm it returns the
constant $\omega_0$.
\end{exercise}

\begin{exercise}
\label{ch25:ex-energy}
Prove that the Hilbert transform preserves energy, $\norm{\hat x}=\norm{x}$, using
Parseval's theorem and the fact that $\abs{-\jj\operatorname{sgn}(\omega)}=1$ for
$\omega\neq0$. What is the one signal class for which energy is \emph{not} preserved, and
why?
\end{exercise}

\begin{exercise}
\label{ch25:ex-orthogonal}
Show that, for a signal with no DC, $x(t)$ and $\hat x(t)$ are orthogonal:
$\int x(t)\hat x(t)\,\dd t = 0$. (Hint: write the inner product in the frequency domain
with Parseval; the multiplier $-\jj\operatorname{sgn}(\omega)$ makes the integrand odd.)
Interpret geometrically in terms of the rotating arrow.
\end{exercise}

\begin{exercise}
\label{ch25:ex-fir}
The ideal sampled Hilbert FIR has taps $h[n]=2/(\pi n)$ for odd $n$ and $0$ for even $n$.
(a) Why are the even taps zero? Relate it to the symmetry of the kernel $1/(\pi t)$.
(b) Truncate to $\abs{n}\le 7$ and apply a Hamming window; compute the resulting taps. (c) Why
must you delay the real signal path by $(\text{length}-1)/2$ samples before forming
$x_a=x_{\text{delayed}}+\jj(h\conv x)$?
\end{exercise}

\begin{exercise}
\label{ch25:ex-code}
Run Listing~\ref{ch25:lst-py}. (a) Plot the envelope and the instantaneous frequency for
the chirp; confirm the frequency sweeps linearly. (b) Add Gaussian noise of standard
deviation $0.05$ to $x$ and recompute -- describe what happens to the instantaneous
frequency and explain why. (c) Smooth the analytic signal (or the phase) before
differentiating and show the improvement.
\end{exercise}

\begin{exercise}
\label{ch25:ex-bedrosian}
Construct a two-tone signal $x(t)=\cos(\omega_1 t)+\cos(\omega_2 t)$ with $\omega_1,
\omega_2$ close. Compute its analytic-signal envelope analytically and show it is
$2\abs{\cos\!\big(\tfrac{\omega_2-\omega_1}{2}t\big)}$ (a beat). Explain why this is a
\emph{correct} envelope yet a \emph{misleading} description if you imagined ``one tone with
slowly varying amplitude,'' and connect this to the Bedrosian warning.
\end{exercise}

\begin{exercise}
\label{ch25:ex-emg}
\textbf{(Toward the wrist tracker.)} You have a $1$\,kHz surface-EMG channel. Design an
envelope follower two ways: (i) Hilbert envelope via a $63$-tap FIR; (ii)
rectify-and-smooth with a one-pole filter at $10$\,Hz cutoff. For each, state the latency,
the multiply-adds per sample, and one accuracy advantage. Then explain, in two sentences,
why neither envelope can tell you \emph{which} finger moved, and what additional sensing
would be required.
\end{exercise}
